\chapter{Conclusions and Future Work}

% =====================================================================
% Cierra el contrato abierto en el Capítulo 1. Sin datos nuevos: todo lo que
% aquí se afirme debe estar respaldado por evidencia ya presentada.
% =====================================================================

\section{Summary of Contributions}
\emph{Recapitular las contribuciones enunciadas en la introducción, ahora
respaldadas por lo construido y lo medido. Una por una, en el mismo orden en que
se enunciaron.}

\section{Answers to the Research Questions}
\emph{Responder explícitamente a cada pregunta de investigación del Capítulo 1,
remitiendo a la evidencia concreta del Capítulo 6 que sostiene cada respuesta.
Sección corta y directa: es la que el tribunal lee con más atención.}

\section{Limitations of the Study}
\emph{Recoger las limitaciones de alcance del trabajo, distinguiéndolas de las
amenazas a la validez ya declaradas en el Capítulo 6: aquí se trata de qué queda
fuera del estudio, no de qué podría sesgar sus resultados.}

\section{Future Work}
\emph{Proponer líneas de continuación derivadas de las limitaciones anteriores,
priorizadas por su valor y no como lista exhaustiva. Cada línea debe conectarse
con una limitación concreta enunciada previamente.}
